\def\topic{Descriptors}
\ifdefined\MyHeaderIncluded
\else
\def\qwertyflag{true}

\ifdefined\MyHeaderIncluded

\else
  \newcommand{\foobar}{FooBar}
  \def\MyHeaderIncluded{true}
% if is till the end of the file

\documentclass[10pt,a4paper]{article}
\usepackage[utf8]{inputenc}
%\usepackage[german]{babel}
%\usepackage[cyrillic]{babel}
\usepackage[T1]{fontenc}
%\usepackage[T2A]{fontenc}
%\usepackage[english, russian]{babel}
\usepackage{amsmath}
\usepackage{amsfonts}
\usepackage{amssymb}
\usepackage{amsthm}
\usepackage{mathtools} %coloneqq

\usepackage{bussproofs} %

%\renewcommand{\thesubsection}{\arabic{subsection}}
%\newtheorem{thm}{Theorem}[section]
%\renewcommand{\thethm}{\arabic{subsection}.\arabic{thm}}

\usepackage{thmtools}
\usepackage{hyperref}
\hypersetup{
    colorlinks=true,
    linkcolor=blue,
    filecolor=magenta,      
    urlcolor=cyan,
    pdftitle={Overleaf Example},
    %pdfpagemode=FullScreen,
    }
%\usepackage{stix-type1}%\turnediota
\theoremstyle{definition}

\newtheorem{constructor}{Constructor}[section]
\newtheorem{iaxiom}{Axiom}
\newtheorem{axiom}{Axiom}
\newtheorem*{axiom*}{Axiom}
\newtheorem{definition}{Definition}[section]
\newtheorem*{definition*}{Definition}
\newtheorem{remark}{Remark}[section]
\newtheorem*{remark*}{Remark}
\newtheorem{recall}{Recall}[section]
\newtheorem*{recall*}{Recall}

\newtheorem{lemma}{Lemma}[section]
\newtheorem{theorem}{Theorem}[section]
\newtheorem*{theorem*}{Theorem}
\newtheorem{note}{Note}[section]
\usepackage{etoolbox}

%\newcommand*{\@Dashv}[2]{%
%  \reflectbox{$\m@th#1#2$}%
%}
\newcommand*{\separation}[3]{% 1)binder, 2)set, 3)condition
  \{#1\in #2: #3\}
}

% https://www.freecodecamp.org/news/simple-regex-tricks-for-beginners-3acb3fa257cb/
%  \\mathcal\{T\} |---> \topology
%  \\mathcal\{([^\}]*)\}   |--->  $1
\newcommand{\AND}{\ \ \&\ \ }
\newcommand{\any}{\mathrm{\ any}}
\newcommand{\borel}{\mathcal{B}}
\newcommand{\can}{\mathrm{can}}
\newcommand{\canonize}{\mathrm{canonize}}
\newcommand{\closure}{\mbox{cl}}
\newcommand{\complex}{\mathbb{C}}
\newcommand{\converges}{\longrightarrow}
\newcommand{\countable}{\mbox{\ countable}}
\newcommand{\ldefined}{\mbox{\ defined}}
\newcommand{\dom}{\mbox{dom}}
\newcommand{\Dom}{\mbox{Dom}}
\newcommand{\Done}{\noindent \textbf{Done.}\\}
\newcommand{\lexists}{\mbox{\ exists}}
\newcommand{\lforall}{\mbox{\ forall}}
\newcommand{\lunique}{\mbox{\ unique}}
\newcommand{\family}{\mathrm{\ family}}
\newcommand{\Fm}{\mathrm{Fm}}
\newcommand{\FuncSymb}{\mathrm{FuncSymb}}
\newcommand{\function}{\mathrm{\ function}}
\newcommand{\functional}{\mbox{\ functional}}
\newcommand{\IFF}{\ \Longleftrightarrow\ }
\renewcommand{\iff}{\ \longleftrightarrow\ }
\newcommand{\IMPLIES}{\ \Longrightarrow\ }
\renewcommand{\implies}{\ \longrightarrow\ }
\newcommand{\Inf}{\mathrm{Inf}}
\newcommand{\injective}{\mbox{\ injective}}
\newcommand{\injection}{\mbox{\ injection}}
\newcommand{\interior}{\mbox{int}}
\newcommand{\intersection}{\bigcap}
\newcommand{\isempty}{\mathrm{\ empty}}
\newcommand{\LB}{\mathrm{LB}}
\newcommand{\nghbset}{\mathcal{N}}
\newcommand{\nonempty}{\mathrm{\ nonempty}}
\newcommand{\Ord}{\mathrm{Ord}}
\newcommand{\power}{\mathcal{P}}
\newcommand{\pr}{\mathrm{pr}}
\newcommand{\fst}{\mathrm{fst}}
\newcommand{\snd}{\mathrm{snd}}
\newcommand{\PredSymb}{\mathrm{PredSymb}}
\newcommand{\Proof}{\noindent \textbf{Proof.}\\}
\newcommand{\ran}{\mathrm{ran}}
\newcommand{\Ran}{\mathrm{Ran}}
\newcommand{\real}{\mathbb{R}}
\newcommand{\relation}{\mathrm{\ relation}}
\newcommand{\set}{\mathrm{\ set}}
\newcommand{\singleton}{\mathrm{\ singleton}}
\newcommand{\semantics}{\mathrm{semantics}}
\newcommand{\semialgebra}{\mathcal{S}}
\newcommand{\Sup}{\mathrm{Sup}}
\newcommand{\surjective}{\mbox{\ surjective\ on\ }}
\newcommand{\Term}{\mathrm{Term}}
\newcommand{\Type}{\mathrm{Type}}
\newcommand{\topology}{\mathcal{T}}
\newcommand{\total}{\mbox{\ total\ on\ }}
\newcommand{\transitive}{\mathrm{\ transitive}}
\newcommand{\union}{\bigcup}
\newcommand{\lTrue}{\mathrm{True}}
\newcommand{\lFalse}{\mathrm{False}}
\newcommand{\lUncertain}{\mathrm{Uncertain}}
\newcommand{\lif}{\mathrm{if}}
\newcommand{\lthen}{\mathrm{then}}
\newcommand{\lelse}{\mathrm{else}}
\newcommand{\lexception}{\mathrm{exception}}
\newcommand{\lundef}{\mathrm{undef}}
\newcommand{\univocal}{\mathrm{\ univocal}}
\newcommand{\univocalat}{\mathrm{\ univocal\ at\ }}
\newcommand{\UB}{\mathrm{UB}}
\newcommand{\Un}{\mathcal{U}}
\newcommand{\Var}{\mathrm{Var}}

\newcommand{\notvarphi}{\phi}
\renewcommand{\phi}{\varphi}

\newcommand{\comment}[1]{\quad\mbox{(#1)}}
% local commands
%\preto\section{\undef\blipo} %etoolbox
\preto\subsection{\undef\blipo} %etoolbox
% https://tex.stackexchange.com/questions/105684/limit-macro-newcommand-scope-to-section
\usepackage{graphicx}
\usepackage{authblk}

\setlength\parindent{0pt} % whole doc of \noindent
\usepackage{soul}
%\usepackage{mathgl}

\fi % guard
\ifdefined\topic
\else
\def\topic{No Titel}
\fi
\title{\topic}
\author[1]{Georgy Dunaev}
\begin{document}
\maketitle
\fi
\section{Descriptors}
\label{sec:descriptors}

This section explores descriptors whose values need not be determined
uniquely by their defining predicates. Its local semantics distinguishes
absence of a candidate, multiplicity of candidates, and a selected value.
It also allows an uncertain truth value $U$. The purpose is to make those
distinctions visible in notation, especially when a descriptor occurs inside
a formula. The final subsection records a separate proposal for axioms of
set-valued descriptors.

\subsection{Existence, uniqueness, and selection}
A predicate can have no satisfying objects (the \emph{nullary} case), exactly
one (the \emph{singular} case), or several (the \emph{plural} case). These
are properties of its possible witnesses. A choice interpretation adds a
further ingredient: it selects one witness when candidates exist. Several
possible witnesses therefore need not make a term ambiguous once a fixed
choice interpretation has been supplied.

The draft uses ``defined'' in places where ``well-defined under the chosen
interpretation'' is intended. We will keep the distinction explicit:
existence concerns whether a value is available; uniqueness concerns whether
the predicate determines it; selection concerns which available value an
interpretation chooses. Invoking choice requires specifying the collection
on which it operates. In particular, a choice principle for sets does not by
itself specify a selector for every predicate on arbitrary classes.

\subsection{Three conventions for a default value}
The superscript in $\tau^D$ describes the admissible default values when the
predicate has no witness. It does not select among several actual witnesses.
There are three conventions:
\begin{enumerate}
\item $\tau^\emptyclass$ supplies no default value. In the nullary case,
  the value is absent; ``at most one'' must not be confused with existence.
\item $\tau^{\{K\}}$ supplies a specified default $K$. The notation
  $\tau^{\exists!}_K$ emphasizes that the default is uniquely specified.
\item $\tau^\Un$ allows any class as a default, where $\Un$ is the
  conglomerate of all classes. This convention leaves the default
  undetermined until an interpretation selects it.
\end{enumerate}
Here the superscript is semantic notation for a collection of admissible
values. If $K$ is a proper class, $\{K\}$ in this notation is not an
object-theoretic singleton class containing $K$.

Our conventions are:
\begin{itemize}
\item $\tau \ceq \tau^\Un$;
\item $\epsilon \ceq \tau^{\exists!}_\emptyclass = \tau^{\{\emptyclass\}}$;
\item $\varepsilon \ceq \tau^\emptyclass$.
\end{itemize}
The total $\epsilon$-operator uses the empty class when no witness exists;
the partial $\varepsilon$-operator supplies no value in that case. The
operators are intended here for definable conglomerates, namely collections
of classes described by predicates in the chosen language. The precise
language and the interpretation of choice are part of the local setting.

\subsection{Quantifiers and total choice}
The following notation separates existence from at-most-one uniqueness:
\[
((\varepsilon x:\varphi) \lunique) \IFF (!x:\varphi),
\qquad
((\varepsilon x:\varphi) \lexists) \IFF (\exists x:\varphi).
\]
Here $!x:\varphi$ means that there is at most one witness, whereas
$\exists!x:\varphi$ asserts both existence and uniqueness. The displayed
conditions describe the witness collection; they are not claims that a
partial term can be substituted without a definedness condition.

With a fixed total choice interpretation, Hilbert-style quantifier clauses
use the total operator $\epsilon$:
\[
(\exists x:\varphi) :\IFF \varphi[(\epsilon x:\varphi)/x],
\qquad
(\forall x:\varphi) :\IFF \varphi[(\epsilon x:\neg\varphi)/x].
\]
Substitution is capture-avoiding. The existential clause selects a witness
when one exists. The universal clause tests a selected counterexample when
one exists; if there is none, every value satisfies $\varphi$, including the
default. This explanation assumes a nonempty domain and ordinary two-valued
predicates under the fixed selection. Extending these clauses to predicates
that themselves take the value $U$ requires further semantic clauses.

The corresponding notation for universal satisfaction is
\[
((\varepsilon x:\varphi) \lforall)
:\IFF (\neg((\varepsilon x:\neg\varphi) \lexists)).
\]
If a partial descriptor has no determined value, substituting it into a
predicate may instead produce the uncertain value $U$. In the strict
connective convention considered here, any connective with an uncertain
argument has uncertain output, including equivalence. Consequently, an
object-language equivalence involving such terms cannot automatically serve
as a two-valued metalinguistic definition. This is the reason to distinguish
the total choice clauses from the partial-descriptor experiment.

\subsection{A two-candidate example}
Let $a$ and $b$ be distinct sets, and consider
\[
c\ceq (\varepsilon x:\ x\in\{a,b\}).
\]
The predicate does not uniquely determine $c$. With a fixed choice
interpretation, however, $c$ denotes one of the two witnesses, so
\[
(a=c)\lor(b=c)
\]
holds under that interpretation. Its individual disjuncts may vary when the
choice interpretation varies. This distinguishes a statement about all
permitted selections from a statement about one selected value.

The exploratory alternative is to leave $c$ unresolved, treating it as a
descriptor with several candidates rather than as an already selected
witness. Then the meanings of $a=c$, $a\simeq c$, and their disjunctions
must be specified before those formulas can be evaluated. A strict
uncertainty convention need not validate the same formulas as fixed choice.
The original question can therefore be stated precisely: which laws of
substitution and equality express membership among the candidates when no
selection has been made?

Introducing a new constant can keep the descriptor visible in the syntax,
but does not by itself answer this semantic question. Likewise, replacing
ordinary equality by $\simeq$ only helps after its interpretation has been
given. These are open design questions for this local calculus, rather than
counterexamples to the ordinary laws of fixed choice.

\subsection{A finite selection model for comparison}
The two-candidate example admits a precise model that keeps track of
correlations between occurrences of a descriptor. This model is a comparison
with the strict $U$ convention, not a replacement for it.

Let $\Omega=\{\omega_a,\omega_b\}$ be two permitted selections. Interpret
$c$ as $a$ at $\omega_a$ and as $b$ at $\omega_b$, keeping the distinct
parameters $a,b$ fixed. Evaluate formulas classically at each selection,
and use the \emph{same} selection at every occurrence of $c$ in a formula.
For a formula $\psi$, let
\[
[\psi]=\{\omega\in\Omega:\psi\text{ is true at }\omega\}.
\]
These brackets denote a set of selections, not an object-theoretic truth
predicate for the whole language. Define a summary value by
\[
v(\psi)=
\begin{cases}
T,&[\psi]=\Omega,\\
F,&[\psi]=\emptyclass,\\
U,&\emptyclass\subsetneq[\psi]\subsetneq\Omega.
\end{cases}
\]
Thus $U$ means variation among permitted classical selections; it does not
mean failure of classical evaluation at any one selection.

\begin{theorem}[Uncertainty summaries lose correlations]
In this model, $v(a=c)=v(b=c)=U$, but
\[
v\bigl((a=c)\lor(b=c)\bigr)=T.
\]
There is no binary truth function on $\{T,F,U\}$ that computes the summary
of every disjunction from the two summaries alone.
\end{theorem}
\Proof
Put $p\equiv(a=c)$. Its truth set is $\{\omega_a\}$, while the truth
set of $b=c$, equivalently $\neg p$ in this model, is $\{\omega_b\}$.
Both summaries are $U$, and the union of their truth sets is $\Omega$.
On the other hand, $p\lor p$ has the same truth set as $p$, so its summary
is $U$. The same input pair $(U,U)$ would therefore require both outputs
$T$ and $U$, which no truth function can provide.
\Done

The full truth sets do determine the connectives:
\[
[\neg\psi]=\Omega\setminus[\psi],\qquad
[\psi\lor\chi]=[\psi]\cup[\chi],\qquad
[\psi\land\chi]=[\psi]\cap[\chi].
\]
These equations follow by evaluating at an arbitrary selection. They make
precise the information that is discarded when a truth set is replaced by
one of the three labels. In particular, $c=c$ is true at every selection,
even though $c$ has two candidates. Treating its two occurrences as
independent choices would describe a different experiment.

The strict convention introduced earlier instead propagates $U$ through
every connective. It can therefore assign $U$ to the displayed disjunction.
The two answers express different meanings of uncertainty: strict propagation
from an unresolved term, versus invariance across complete selections.
A calculus should say which meaning it implements before adopting either
set of laws.

\subsection{Substitution must transport the selections}
The shared-selection interpretation distinguishes truth throughout one
specified experiment from logical validity under all assignments.
This distinction matters when a rule permits substitution. The following
result does not require finitely many selections: $\Omega$ is a supplied
set of complete assignments. It does not assert that every family of
nonempty candidate sets admits a complete selection, and thus does not
introduce a hidden choice principle.

\begin{theorem}[Transport of truth sets under substitution]
Work with a classical first-order structure whose terms have total values,
and let $\Omega$ be a nonempty set of assignments. For a finite,
sort-correct, capture-avoiding simultaneous substitution $\sigma$, define
\[
F_\sigma(v)(x)=\llbracket x\sigma\rrbracket_v\quad(v\in\Omega),
\qquad \Omega_\sigma=F_\sigma[\Omega].
\]
Here $F_\sigma$ has domain $\Omega$ and codomain $\Omega_\sigma$;
the inverse image below is taken for this restricted map. Writing
$[\phi]_S$ for the assignments in $S$ satisfying $\phi$,
\[
[\phi\sigma]_\Omega
 =F_\sigma^{-1}\bigl([\phi]_{\Omega_\sigma}\bigr).
\]
In particular, $\phi\sigma$ holds throughout $\Omega$ exactly when
$\phi$ holds throughout $\Omega_\sigma$. If
$\Omega_\sigma\subseteq\Omega$, truth throughout $\Omega$ is preserved
by this substitution.
\end{theorem}
\Proof
The ordinary substitution lemma gives
$v\models\phi\sigma$ if and only if $F_\sigma(v)\models\phi$.
This proves the inverse-image equality. Every member of $\Omega_\sigma$
has a preimage by its definition, so universal truth on the two sides is
equivalent. The inclusion in the last claim lets truth on $\Omega$ be
restricted to $\Omega_\sigma$.
\Done

\begin{example}[Universal truth in an experiment is not substitution validity]
Let $D=\{0,1\}$, and let $\Omega$ consist of assignments in which
$x=y$, allowing both common values. Then $x=y$ is true throughout
$\Omega$. Substituting the constants $0$ and $1$ for $x$ and $y$
produces $0=1$, which is false throughout $\Omega$. The transported
assignments have $x=0,y=1$ and lie outside $\Omega$.
There is no conflict with classical substitution of valid formulas:
$x=y$ was valid only on this restricted experiment, not on all assignments.
\end{example}

\begin{proposition}[Marginal equality does not justify correlated replacement]
Equality of the candidate sets of two terms need not permit substitution
of one for the other in shared-selection formulas.
\end{proposition}
\Proof
On $\Omega=\{\omega_0,\omega_1\}$, take $c(\omega_i)=i$ and
$d(\omega_i)=1-i$. Both terms have candidate set $\{0,1\}$.
Nevertheless, the context $z=c$ is true everywhere after substituting
$c$ for $z$, and false everywhere after substituting $d$ for $z$.
Their marginal sets agree, but their joint values with $c$ do not.
\Done

This does not invalidate equality of candidate sets in the alternative
powerset-domain model: there a candidate set is itself the semantic object,
and no further selection correlation is retained. In a shared-selection
model, the sufficient replacement condition is pointwise agreement
$s(v)=t(v)$ at every admitted assignment (together with the ordinary
binder freshness conditions). Equality of their ranges is weaker.

\subsection{Proposed axioms for set-valued descriptors}
This subsection uses descriptors restricted to set witnesses. Read the
following formulas as proposed local axiom schemes. The notation
$\ldefined$ concerns the descriptor's value; the earlier strict propagation
rule for $U$ does not supply a semantics for these schemes on its own.

\begin{axiom}[Choice of an existing set witness]
\[
\exists (x\set):\varphi
\implies
(\varepsilon (x\set): \varphi) \in \cc{x}{\varphi}
\]
\end{axiom}
\begin{axiom}[Existence required for partial choice]
\[
(\varepsilon (x\set): \varphi) \ldefined
\implies
\exists (x\set):\varphi
\]
\end{axiom}
\begin{axiom}[Description of a unique set witness]
\[
(\exists! (x\set):\varphi)
\implies
(\iota (x\set): \varphi) \in \cc{x}{\varphi}
\]
\end{axiom}
\begin{axiom}[Uniqueness required for description]
\[
(\iota (x\set): \varphi) \ldefined
\implies
\exists! (x\set):\varphi
\]
\end{axiom}

A related encoding represents an unsuccessful description by a proper-class
default, so that sethood can detect success. Under that additional convention
one proposes
\[
(\iota (x\set): \varphi) \ldefined
\iff
(\iota (x\set): \varphi) \set.
\]
This is an encoding of definedness for set-valued descriptions, not a claim
that proper classes are intrinsically undefined.

The intended combined specification is
\[
(\iota (x\set): \varphi) \ldefined
\iff
\exists! (x\set):\varphi,
\]
with the witness property
\[
(\iota (x\set): \varphi) \ldefined
\implies
(\iota (x\set): \varphi) \in \cc{x}{\varphi}.
\]
For example, if $x$ is a set satisfying $\varphi$ and there is at most one set
witness, the description is defined:
\[
(x\set)\land\varphi\land (!(x\set):\varphi)
\implies
(\iota (x\set): \varphi) \ldefined.
\]
The free occurrence of $x$ in the premise supplies existence. The sethood
condition makes explicit the witness domain used by this subsection.

\ifdefined\qwertyflag
\undef\qwertyflag
\end{document}
\fi
